\chapter{Metodología}

Este capítulo describe exhaustivamente el diseño experimental empleado para validar la hipótesis central de esta tesis. Se detalla cada componente del \textit{pipeline} computacional: desde la selección y preprocesamiento de los datos de entrada, pasando por la construcción de grafos y la simulación controlada de fallas sensoriales, hasta la optimización bayesiana de hiperparámetros y la evaluación cuantitativa final. El principio rector fue garantizar que todas las familias de métodos ---geométricos, estadísticos y basados en grafos--- compitan bajo condiciones estrictamente equivalentes de sintonización paramétrica, eliminando así el sesgo de benchmarking que permea la literatura actual.

\section{Arquitectura General del Pipeline}

\begin{figure}[ht]
    \centering
    \begin{tikzpicture}[
        box/.style={rectangle, rounded corners, draw=blue!70!black, fill=blue!5, thick, minimum height=1cm, text centered, font=\small},
        arrow/.style={-latex, thick, blue!70!black}
    ]
        % Nodes
        \node (data) [box, align=center] {Carga y preprocesamiento\\(MNE-Python)};
        \node (mask) [box, right=of data, align=center] {Máscara de Degradación\\(Random vs Nearby)};
        \node (opt) [box, right=of mask, align=center, fill=orange!10, draw=orange!80!black] {Optimización Bayesiana\\(Optuna + TPE/NSGA-II)};
        \node (buffer) [box, below=1cm of opt, align=center] {Buffer Cronológico (2s)\\(Prevención de Leakage)};
        \node (eval) [box, left=of buffer, align=center] {Evaluación Multi-Métrica\\(MAE, DTW, LSD)};
        
        % Arrows
        \draw [arrow] (data) -- (mask);
        \draw [arrow] (mask) -- (opt);
        \draw [arrow] (opt) -- (buffer);
        \draw [arrow] (buffer) -- (eval);
    \end{tikzpicture}
    \caption{Representación esquemática del pipeline metodológico de optimización y evaluación propuesto.}
    \label{fig:methodology_pipeline}
\end{figure}

El flujo de trabajo se organiza como una secuencia determinista de seis etapas. Primero se realiza la carga y estandarización de los datasets heterogéneos mediante MNE-Python, seguida de la construcción del grafo sobre las coordenadas 3D de los electrodos y/o la covarianza temporal de la señal. A continuación se aplica la simulación de canales faltantes bajo protocolos espaciales controlados (\textit{random} y \textit{nearby}), tras lo cual se ejecuta la interpolación y reconstrucción mediante el abanico completo de métodos implementados. Posteriormente, los hiperparámetros de cada método son sintonizados mediante Optuna con muestreo TPE y NSGA-II multi-objetivo. Finalmente, se realiza la evaluación cuantitativa con seis métricas complementarias y la consolidación estadística de los resultados.

Cada etapa genera artefactos intermedios serializados (archivos CSV, JSON y matrices NumPy), lo que permite desacoplar el cómputo intensivo de la generación de figuras y tablas finales. Esta modularidad garantiza la trazabilidad completa desde el dato crudo hasta cada cifra reportada en el manuscrito.

Como capa final de control experimental, la cobertura de configuraciones se gobernó mediante un registro exhaustivo que define explícitamente el espacio combinatorio de datasets, constructores de grafo, escenarios de pérdida, semillas, hiperparámetros y familias de métodos. De esta forma, cada resultado reportado mantiene una trazabilidad completa desde la configuración inicial hasta los artefactos generados (tablas de estadísticas, metadatos y figuras).

La implementación computacional se desarrolló estructurando el código en módulos lógicos que separan responsabilidades. Un primer módulo se encarga de la carga unificada de los tres datasets. Otro módulo agrupa los constructores de grafos espaciales y temporales. Las funciones de interpolación reúnen tanto los algoritmos GSP como los \textit{baselines} geométricos e ICA. Un módulo de evaluación centraliza el cálculo de todas las métricas de rendimiento, y finalmente los scripts orquestadores gestionan la simulación, la búsqueda bayesiana de hiperparámetros y la visualización de resultados.

\section{Bases de Datos Experimentales}

La validez externa de cualquier estudio de interpolación depende críticamente de la diversidad de los datos de prueba. Se seleccionaron tres datasets públicos que cubren un amplio espectro de densidades espaciales, frecuencias de muestreo y paradigmas cognitivos.

\subsection{PhysioNet EEG Motor Movement/Imagery Dataset (EEGMMIDB)}

Este dataset, alojado en la base de datos PhysioNet, comprende registros de 64 canales EEG adquiridos a 160~Hz bajo el sistema internacional 10-10. Los sujetos realizaron tareas de ejecución motora real (apertura y cierre de puños) e imagería motora (imaginación del mismo movimiento sin ejecución). La alta densidad de electrodos genera una malla espacial densa con alta correlación inter-canal, lo que constituye el escenario más favorable para los métodos geométricos.

La carga y estandarización se realizan mediante las funciones provistas por MNE-Python, aplicando la nomenclatura estándar de canales y asignando el montaje 10-05 para obtener coordenadas 3D.

\subsection{BCI Competition IV Dataset 2a}

El dataset BCI Competition IV 2a \cite{tangermann2012bci} es un estándar en la comunidad de interfaces cerebro-computador. Contiene registros de 22 canales EEG (montaje 10-20 extendido) adquiridos a 250~Hz de 9 sujetos realizando imaginería motora de 4 clases: mano izquierda, mano derecha, ambos pies y lengua.

La baja densidad espacial de este dataset (22 canales frente a los 64 de PhysioNet) representa un desafío extremo para la interpolación. Con solo 22 nodos en el grafo, la pérdida de incluso 2--3 canales puede dejar regiones del cráneo sin cobertura sensora alguna, exponiendo las debilidades inherentes de los métodos puramente geométricos.

Los archivos en formato GDF requieren descarga manual desde el sitio oficial de la competencia. La carga se realiza mediante las funciones de lectura de MNE-Python, con un mapeo explícito de nombres de canal al estándar 10-05.

\subsection{MNE Sample Dataset}

El MNE Sample Dataset \cite{gramfort2013meg} es un registro multimodal (MEG y EEG) de 60 canales EEG a 600~Hz, capturando potenciales evocados auditivos y visuales. Este dataset fue seleccionado específicamente por la presencia de componentes ERP bien documentados (N100, P300), cuya preservación morfológica y temporal constituye el criterio más exigente de evaluación clínica.

La carga se realiza mediante MNE-Python, seleccionando exclusivamente los canales EEG y extrayendo las posiciones 3D del montaje asociado.

\subsection{Resumen Comparativo de Datasets}

\begin{table}[ht]
\centering
\caption{Características de los tres datasets experimentales.}
\label{tab:datasets}
\begin{tabular}{lccc}
\toprule
\textbf{Característica} & \textbf{PhysioNet} & \textbf{BCI IV 2a} & \textbf{MNE Sample} \\
\midrule
Canales EEG & 64 & 22 & 60 \\
Frecuencia (Hz) & 160 & 250 & 600 \\
Paradigma & Motor MI & Motor MI 4-cl & Evocados A/V \\
Montaje & 10-10 & 10-20 ext. & Personalizado \\
Densidad espacial & Alta & Baja & Alta \\
\bottomrule
\end{tabular}
\end{table}

\section{Construcción de Grafos}

La topología del grafo es un factor determinante en el rendimiento de cualquier algoritmo GSP. En esta tesis se evaluaron \ResNGraphMethods~constructores de grafo organizados en tres categorías.

\subsection{Grafos de Vecindad Geométrica}

Estos grafos conectan electrodos basándose exclusivamente en la proximidad euclidiana de sus coordenadas 3D en el cuero cabelludo. Se evaluaron cinco variantes.

El constructor más elemental es el \textbf{$k$-Nearest Neighbors} ($k$-NN), que conecta cada electrodo con sus $k$ vecinos más cercanos y aplica simetrización posterior por máximo para garantizar un grafo no dirigido. Su extensión natural es el \textbf{$k$-NN Gaussiano} ($k$-NNG), que asigna pesos gaussianos a las aristas:
\begin{equation}
    W_{ij} = \exp\left(-\frac{d(i,j)^2}{2\sigma^2}\right) \cdot \mathbb{1}_{[j \in \mathcal{N}_k(i)]}
\end{equation}
donde $\sigma$ controla el ancho del kernel y $\mathcal{N}_k(i)$ denota el conjunto de $k$ vecinos de $i$. Inspirado en Tamaru et al.~\cite{tamaru_optimizing_2024}, el \textbf{$k$-NN Variable Gaussiano} adapta $k_i$ por nodo según la densidad local de sensores:
\begin{equation}
    k_i = \text{clip}\left(\text{round}\left(k_{\text{base}} \cdot \left(\frac{\tilde{d}}{d_i}\right)^\alpha\right), k_{\min}, k_{\max}\right)
\end{equation}
donde $\tilde{d}$ es la mediana de las distancias al $k$-ésimo vecino y $d_i$ la distancia local del nodo $i$. Esta variante resulta especialmente útil en montajes con densidad de electrodos no uniforme.

Dos constructores adicionales completan la familia geométrica. El grafo de \textbf{bola-$\epsilon$} conecta todos los pares $(i,j)$ cuya distancia satisface $d(i,j) < \epsilon$, mientras que el \textbf{árbol de expansión mínima} (MST, \textit{Minimum Spanning Tree}) garantiza la conectividad mínima del grafo.

\subsection{Grafos Densos / Globales}

Se evaluaron también dos topologías densas: un grafo \textbf{completamente conectado con pesos de distancia inversa}, donde $W_{ij} = 1/d(i,j)$ para todo par de electrodos, y un grafo \textbf{gaussiano global} sin restricción de vecindad.

\subsection{Grafos Adaptativos e Inferencia Basada en Datos}

Estos métodos aprenden la topología directamente desde los datos, superando las limitaciones del modelo esférico. El constructor \textbf{NNK} (\textit{Non-Negative Kernel Regression})~\cite{shekkizhar2020} resuelve un problema de mínimos cuadrados no negativos local para cada nodo, obteniendo pesos de arista esparcidos que reflejan genuinamente la estructura geométrica local. El método \textbf{AEW} (\textit{Adaptive Edge Weighting})~\cite{karasuyama_adaptive_2017} optimiza los pesos de las aristas minimizando un error de reconstrucción local bajo una parametrización por funciones de similitud; cuando el optimizador no converge, un mecanismo de respaldo combina kernel gaussiano espacial con correlación temporal. Finalmente, el modelo de \textbf{aprendizaje de grafos de Kalofolias}~\cite{kalofolias_how_2016} infiere la topología con regularización \textit{log-degree}:
\begin{equation}
    \min_{\mathbf{W} \geq 0} \langle \mathbf{Z}, \mathbf{W} \rangle - a \sum_i \log(d_i) + \frac{b}{2} \|\mathbf{W}\|_F^2
\end{equation}
donde $\mathbf{Z}$ es la matriz de distancias al cuadrado entre señales, y los parámetros $a, b$ equilibran esparsidad y conectividad. La solución se obtiene por descenso de gradiente proyectado con restricciones de simetría y no-negatividad.

Tras una fase exploratoria inicial en la que se evaluaron todos los constructores sobre los tres datasets, se seleccionó el $k$-NN Gaussiano como constructor primario para la fase de optimización final, dada su estabilidad numérica (ausencia de grafos disconexos), su eficiencia computacional y su desempeño competitivo sostenido a lo largo de todos los escenarios.

\section{Protocolo de Simulación de Canales Faltantes}

Para evaluar la robustez bajo condiciones realistas, se diseñaron dos modos de falla espacial que capturan los escenarios más frecuentes en la práctica clínica.

\subsection{Modo \textit{Random} (Fallas Independientes)}

Los canales se seleccionan para remoción mediante un muestreo uniforme discreto sin reemplazo. Este modo simula fallas estocásticas e independientes, como el deterioro gradual del gel conductor en electrodos arbitrarios o el descarte manual de canales ruidosos durante el preprocesamiento.

\subsection{Modo \textit{Nearby} (Fallas Contiguas)}

Se selecciona un electrodo semilla aleatorio y se remueven sus $m$ vecinos más cercanos (por distancia euclidiana 3D). Este modo simula fallas correlacionadas espacialmente de alta adversidad: desprendimiento de un parche del gorro EEG, falla de un módulo de amplificación local, o saturación de toda una región por un artefacto de movimiento.

\subsection{Niveles de Degradación}

Cada modo se combina con dos esquemas de severidad: por proporción (\textit{ratio}), removiendo el 10\%, 20\%, 30\% o 40\% de los canales totales, y por conteo discreto (\textit{count}), removiendo 1, 2 o 3 canales específicos. La combinación de 2 modos $\times$ 7 niveles produce 14 escenarios experimentales por dataset, resultando en 42 contextos experimentales únicos.

\section{Métodos de Interpolación y Reconstrucción}

El framework implementa un total de 18 algoritmos de interpolación organizados en tres familias.

\subsection{Familia 1: Baselines Geométricos y Estadísticos}

Estos métodos no utilizan información de grafos y sirven como referencia inferior.

\begin{itemize}
    \item \textbf{Spherical Spline} \cite{perrin_spherical_1989}: Estándar \textit{de facto} en herramientas clínicas como MNE-Python y EEGLAB. Modela el cuero cabelludo como una esfera e interpola usando polinomios de Legendre. El hiperparámetro de regularización $\varepsilon$ controla la suavidad.

    \item \textbf{RBF Interpolation (Thin-Plate Spline y Multiquadric)} \cite{jager_reconstruction_2016}: Funciones de base radial con kernel \textit{Thin-Plate} o \textit{Multiquadric}. El parámetro de suavizado controla el \textit{trade-off} entre fidelidad a los datos e interpolación suave.

    \item \textbf{ICA (Independent Component Analysis)}: Basado en el algoritmo FastICA. Los canales faltantes se inicializan con la media del canal y se reconstruyen desde la proyección inversa de las componentes independientes.

    \item \textbf{IDW (Inverse Distance Weighting)}: Promedio ponderado por la inversa de la distancia euclidiana elevada a una potencia $p$.

    \item \textbf{Nearest Neighbor / Mean}: Baselines triviales para establecer el piso de rendimiento absoluto.
\end{itemize}

\subsection{Familia 2: Métodos GSP Espaciales (Instantáneos)}

Estos métodos operan muestra por muestra ($t$ fijo), resolviendo un problema de optimización espacial independiente en cada instante temporal.

\begin{itemize}
    \item \textbf{GSP puro (Laplacian Solving)}: Resuelve el sistema $\mathbf{L}_{\mathcal{U}\mathcal{U}} \hat{\mathbf{x}}_\mathcal{U} = -\mathbf{L}_{\mathcal{U}\mathcal{K}} \mathbf{x}_\mathcal{K}$ directamente mediante la factorización del subsistema lineal correspondiente.

    \item \textbf{Tikhonov en Grafo}: Resuelve $(\mathbf{M} + \alpha \mathbf{L})\hat{\mathbf{x}} = \mathbf{M}\mathbf{y}$, donde $\mathbf{M}$ es la máscara diagonal de observación y $\alpha$ el peso de regularización.

    \item \textbf{BGSRP} \cite{zhang_recovery_2024}: Reconstrucción en RKHS bandlimited. Utiliza los primeros $B$ autovectores del Laplaciano (excluyendo la componente DC) para construir un kernel $\mathbf{K}$ y resolver con regularización $\gamma$.

    \item \textbf{Graph Smoothing}: Difusión iterativa sobre la matriz de transición normalizada del grafo.

    \item \textbf{Graph TV}: Variación total en grafo via IRLS (\textit{Iteratively Reweighted Least Squares}), promoviendo señales con pocas discontinuidades abruptas.

    \item \textbf{Sobolev / Puy}: Variantes de regularización con operadores diferenciales de alto orden sobre el grafo.
\end{itemize}

\subsection{Familia 3: Métodos con Regularización Temporal}

Estos métodos operan sobre ventanas temporales completas $\mathbf{X} \in \mathbb{R}^{N \times T}$, optimizando conjuntamente los gradientes espaciales y temporales. Esta es la familia que constituye la contribución analítica central de esta tesis.

\begin{itemize}
    \item \textbf{TRSS} (Temporal Regularized Sobolev Smoothness) \cite{giraldo2022}: Minimiza la función de costo espacio-temporal completa. La implementación utiliza descenso de gradiente con paso adaptativo basado en la cota de Lipschitz:
    \begin{equation}
        L = 2\left(1 + \alpha \|\mathbf{L}\|_2 + \beta \|\mathbf{L}_t\|_2 \|\mathbf{L}\|_2\right)
    \end{equation}
    donde $\mathbf{L}_t = \mathbf{D}_t^\top \mathbf{D}_t$ es el Laplaciano temporal (operador de diferencias finitas). Las observaciones conocidas se fijan en cada iteración (\textit{hard constraint}).

    \item \textbf{Graph-Time Tikhonov}: Combina regularización Tikhonov espacial con suavizado temporal bilateral en dos etapas secuenciales.

    \item \textbf{Sobolev Temporal / Temporal Laplacian / Heat Diffusion Temporal}: Variantes que modifican la formulación del operador temporal.
\end{itemize}

\section{Optimización de Hiperparámetros con Optuna}

La contribución metodológica más importante de este trabajo es someter \textit{todos} los métodos ---incluyendo los \textit{baselines} clásicos--- a una optimización de hiperparámetros rigurosa. Se utilizó el \textit{framework} \textbf{Optuna}~\cite{akiba2019optuna} por su eficiencia en espacios de búsqueda continuos y mixtos.

\subsection{Reducción del Espacio de Métodos}

Aunque la Sección~3.5 describió un total de 18 algoritmos de interpolación, la fase de optimización bayesiana no se aplicó de forma indiscriminada a todos ellos. Se realizó primero una fase exploratoria inicial con hiperparámetros por defecto sobre todos los datasets y escenarios de pérdida, cuyo propósito fue identificar las familias competitivas y descartar aquellas con rendimiento consistentemente inferior o con inestabilidad numérica recurrente (por ejemplo, métodos puramente laplacianos sin regularización, o BGSRP bajo pérdida contigua severa). A partir de esta criba, la optimización fina con Optuna se concentró en los métodos con mayor potencial: \textbf{Tikhonov en grafo} y \textbf{TRSS} como representantes de la familia GSP, junto con los tres \textit{baselines} clásicos (Spherical Spline, RBF y ICA) para mantener una comparación justa. Análogamente, del espacio de constructores de grafo se retuvo únicamente el $k$-NN Gaussiano, por las razones expuestas en la Sección~3.3.

\subsection{Estrategia de Muestreo}

Para la optimización mono-objetivo (minimizar MAE), se empleó el muestreador \textbf{TPE} (\textit{Tree-structured Parzen Estimator}), que modela la distribución condicional de los hiperparámetros dado el rendimiento observado, concentrando la exploración en regiones prometedoras.

Para la optimización multi-objetivo (minimizar simultáneamente MAE y LSD), se empleó el muestreador \textbf{NSGA-II} (\textit{Non-dominated Sorting Genetic Algorithm II}), que construye un frente de Pareto de soluciones no dominadas. La selección del mejor trial del frente se realiza por distancia euclidiana normalizada al origen.

\subsection{Espacios de Búsqueda por Familia}

Los rangos de búsqueda fueron definidos para cubrir las regiones de interés sin incluir configuraciones degeneradas:

\begin{table}[ht]
\centering
\caption{Espacios de búsqueda de hiperparámetros por método.}
\label{tab:hyperparams}
\begin{tabular}{lll}
\toprule
\textbf{Método} & \textbf{Parámetro} & \textbf{Rango} \\
\midrule
Spherical Spline & $\varepsilon$ (regularización) & $[10^{-8}, 10^{-2}]$ (log) \\
RBF (TPS) & suavizado & $[0.0, 5.0]$ \\
ICA & $n$ (componentes) & $[1, N_{\text{ch}}]$ (entero) \\
\midrule
Tikhonov & $\alpha$ & $[10^{-4}, 10^{2}]$ (log) \\
\midrule
TRSS & $\alpha$, $\beta$ & $[10^{-3}, 10]$, $[10^{-4}, 1]$ \\
TRSS & iteraciones, tasa de aprendizaje & $[50, 500]$, $[10^{-4}, 0.1]$ \\
\midrule
Grafo ($k$-NNG) & $k$, $\sigma$ & $[2, N/2]$, $[0.01, 5]$ \\
\bottomrule
\end{tabular}
\end{table}

\subsection{Protocolo de Ejecución}

Para cada combinación (dataset $\times$ escenario de falla $\times$ método), Optuna fue ejecutado siguiendo una política de muestreo estratificada. Los métodos GSP (Tikhonov y TRSS) emplearon el muestreador multi-objetivo NSGA-II con 30~\textit{trials} por combinación, optimizando simultáneamente MAE y LSD. Dentro de cada \textit{trial} de TRSS, el algoritmo iterativo de descenso de gradiente ejecuta internamente hasta 80~iteraciones para converger a la solución del problema de optimización espacio-temporal definido en la Ecuación~\ref{eq:trss_obj}. Los \textit{baselines} geométricos (Spherical Spline, RBF e ICA) se optimizaron con 10~\textit{trials} por combinación. En total, el protocolo ejecutó del orden de miles de evaluaciones de interpolación al cruzar los 3~datasets, 14~escenarios de pérdida y los métodos optimizados. La selección del mejor \textit{trial} desde el frente de Pareto resultante se realizó por distancia euclidiana normalizada al origen, y los hiperparámetros óptimos junto con las métricas resultantes se consolidaron en tablas estructuradas para su análisis posterior.

\subsection{Validación Temporal y Prevención de Fuga de Datos}

Un aspecto crítico en la optimización de hiperparámetros para señales temporales es evitar que el modelo acceda a los mismos datos durante la sintonización y la evaluación final. Para mitigar este riesgo se implementó un esquema de validación temporal estricto. Cada segmento de señal se divide cronológicamente en dos bloques: los primeros 70\% de las muestras constituyen el conjunto de optimización, mientras que el bloque restante se reserva como conjunto de prueba (\textit{Test Set}). Entre ambos bloques se descarta un \textit{buffer} de 2~segundos ($2 \times f_s$ muestras) para contrarrestar la fuerte autocorrelación temporal inherente a señales EEG filtradas por banda pasante (0.5--40~Hz); esta zona de separación asegura que las muestras de prueba no sean trivialmente predecibles desde la frontera de optimización. Optuna evalúa la función objetivo exclusivamente sobre el conjunto de optimización, y una vez seleccionados los mejores parámetros, todas las métricas reportadas en los capítulos de Resultados y Discusión se calculan únicamente sobre el conjunto de prueba.

Este protocolo garantiza que los resultados reflejen la capacidad genuina de los algoritmos para reconstruir información neurofisiológica futura, simulando una aplicación en tiempo real donde los parámetros óptimos del pasado deben generalizar a muestras nuevas.

\section{Evaluación Cuantitativa}

El rendimiento de cada método se cuantifica mediante seis métricas complementarias que capturan dimensiones ortogonales del error de reconstrucción.

En el dominio de la amplitud se emplean tres indicadores. El \textbf{MAE} (Mean Absolute Error), definido como $\text{MAE} = \frac{1}{T} \sum_{t=1}^{T} |x(t) - \hat{x}(t)|$, mide la desviación absoluta promedio y resulta robusto ante \textit{outliers}. El \textbf{RMSE} (Root Mean Square Error), $\text{RMSE} = \sqrt{\frac{1}{T} \sum (x(t) - \hat{x}(t))^2}$, penaliza severamente errores grandes aislados, funcionando como detector de inestabilidades numéricas. Finalmente, el \textbf{SNR} (Signal-to-Noise Ratio), $\text{SNR} = 10 \log_{10}\left(\frac{\text{Var}(x)}{\text{Var}(x - \hat{x})}\right)$~dB, proporciona una medida global fácilmente interpretable entre escalas de sensores distintas.

Para evaluar la fidelidad temporal se utiliza el \textbf{DTW} (Dynamic Time Warping), que calcula la alineación temporal óptima entre dos secuencias minimizando la distancia euclidiana acumulada. Esta métrica es esencial para la evaluación de ERPs, ya que tolera pequeños desfases de latencia sin penalizarlos como error de amplitud. Se emplea programación dinámica con ventana de Sakoe-Chiba para acotar la complejidad computacional.

En el dominio espectral se emplean dos métricas adicionales. La \textbf{LSD} (Log Spectral Distance), definida como $\text{LSD} = \sqrt{\frac{1}{K} \sum_{k} (\log |X(k)| - \log |\hat{X}(k)|)^2}$, cuantifica la preservación del contenido frecuencial en las bandas fisiológicas, donde $X(k)$ y $\hat{X}(k)$ son las magnitudes de la FFT de la señal original y reconstruida respectivamente. La \textbf{Coherencia} (Magnitude-Squared Coherence), $C_{xy}(f) = \frac{|S_{xy}(f)|^2}{S_{xx}(f) S_{yy}(f)}$, complementa el análisis indicando la preservación de la correlación de fase en cada frecuencia, con valores cercanos a 1 representando una reconstrucción espectral ideal.

Todas las métricas se promedian sobre los canales reconstruidos. La combinación de estas seis dimensiones impide que un método oculte sus debilidades: un MAE bajo obtenido por sobre-suavizado se evidenciará como un DTW y LSD elevados.

\section{Consolidación Estadística y Generación de Artefactos}

Los resultados individuales se agregan mediante tablas pivote que resumen el rendimiento óptimo por familia de método, dataset y escenario. A partir de estas tablas se construyen rankings ordinales por escenario que permiten identificar la familia dominante en cada contexto experimental. Adicionalmente, se aplican pruebas de significancia estadística no paramétricas a nivel de \textit{trial} mediante Mann-Whitney U y delta de Cliff para validar que las diferencias observadas no sean producto del azar.

La visualización de los resultados se apoya en diversos formatos complementarios: mapas de calor (\textit{heatmaps}) que cruzan métodos con escenarios para la métrica RMSE, gráficos radar multi-métrica que permiten comparar simultáneamente las seis dimensiones normalizadas, curvas de degradación que muestran la caída del SNR en función de la proporción de canales faltantes, frentes de Pareto que evidencian el \textit{trade-off} entre DTW y MAE, y series de tiempo de ERPs junto con sus correspondientes PSD para la comparación cualitativa visual de la reconstrucción.

\section{Stack Tecnológico e Infraestructura}

La totalidad del \textit{pipeline} se implementó en \textbf{Python~3.10+}, aprovechando el ecosistema científico abierto. La carga, preprocesamiento y manipulación de datos EEG se realiza con \textbf{MNE-Python}~\cite{gramfort2013meg}. El álgebra lineal densa, las descomposiciones espectrales y las funciones de interpolación se apoyan en \textbf{NumPy} y \textbf{SciPy}, mientras que \textbf{scikit-learn} provee la construcción de grafos $k$-NN y el algoritmo FastICA. La optimización bayesiana de hiperparámetros se gestiona mediante \textbf{Optuna}~\cite{akiba2019optuna}, y la consolidación tabular junto con la generación de figuras se realiza con \textbf{Pandas} y \textbf{Matplotlib}. Para el cálculo eficiente de DTW se emplea la biblioteca \textbf{dtaidistance}, que ofrece una implementación optimizada en C.

Todas las operaciones de álgebra matricial asociadas a la inversión de los sistemas lineales GSP (Tikhonov, BGSRP) se implementaron con rutinas vectorizadas propias para maximizar la reproducibilidad y el control sobre las tolerancias numéricas. Para BGSRP se implementó además una pseudo-inversa con tolerancia numérica equivalente a la utilizada en implementaciones de referencia.

\section{Reproducibilidad y Transparencia}

La reproducibilidad total del experimento se garantiza mediante diversas medidas complementarias. Todas las operaciones estocásticas ---muestreo de máscaras, inicialización de ICA y muestreo de Optuna--- utilizan semillas deterministas, asegurando resultados idénticos entre ejecuciones. Los archivos de resultados, configuraciones experimentales y bitácoras de ejecución se almacenan bajo control de versiones en un repositorio Git. Cada fase del \textit{pipeline} puede re-ejecutarse independientemente, lo que permite reproducir de forma aislada la optimización, la evaluación o la generación de figuras. Además, un módulo dedicado captura y serializa cualquier advertencia numérica ---como problemas de convergencia de ICA o singularidad del Laplaciano--- para facilitar la auditoría posterior de la estabilidad numérica del sistema.
